\documentclass{article}
\usepackage[UTF8]{ctex}
\usepackage{amsmath,amssymb,amsfonts}
\usepackage{enumitem}
\usepackage{geometry}
\geometry{a4paper, margin=1in}
% === 错题本排版优化（不参与题目标识） ===
\setlength{\parindent}{0pt}
\setlength{\parskip}{0.35em plus 0.10em minus 0.08em}
\setlength{\emergencystretch}{3em}
\setlist{itemsep=0.25em, topsep=0.40em, parsep=0pt}
\allowdisplaybreaks[2]
% === 错题本排版优化结束 ===
\begin{document}
\section*{中国科学技术大学2024年数学分析真题}
\section*{104．中国科学技术大学2024年数学分析试卷}

\subsection{题目1}
1．（15 分）设 \(f=f(u, v, w)\) 是 \(\mathbb{R}^{3}\) 上定义的光滑函数，求 \(g(x, y, z)=f\left(x+y, x y, \frac{x}{y}\right)\) 的偏导数\(\frac{\partial g}{\partial x}, \frac{\partial^{2} g}{\partial x^{2}}\) ．

\subsection{题目2}
2．（15 分）设 \(a>0\) ，求旋轮线一拱 \(\left\{\begin{array}{l}x=a(t-\sin t) ; \\ y=a(1-\cos t) .\end{array} 0 \leq t \leq 2 \pi\right.\) 的弧长．

\subsection{题目3}
3．（15 分）计算积分

\[
\iint_{S}\left(x^{2}+y\right) \mathrm{d} y \mathrm{~d} z-y \mathrm{~d} z \mathrm{~d} x+\left(x^{2}+y^{2}\right) \mathrm{d} x \mathrm{~d} y .
\]

其中 \(S\) 为旋转抛物面 \(z=x^{2}+y^{2}\) 在 \(z \leq 1\) 的部分，并取曲面外侧为正向。

\subsection{题目4}
4．（15 分）计算极限 \(\lim _{n \rightarrow \infty} n\left(\frac{1}{n^{2}+1^{2}}+\frac{1}{n^{2}+2^{2}}+\cdots+\frac{1}{n^{2}+n^{2}}\right)\) ．

\subsection{题目5}
5．（15 分）确定函数 \(I(\alpha)=\int_{0}^{+\infty} \frac{\ln \left(1+x^{3}\right)}{x^{\alpha}} \mathrm{d} x\) 的定义域，并证明 \(I\) 在定义域上是连续函数。

\subsection{题目6}
6．（15 分）设 \(a>0, x_{0}>0\) 而 \(x_{n}\) 由递推公式 \(x_{n}=\frac{1}{2}\left(x_{n-1}+\frac{a}{x_{n-1}}\right)(n \geq 1)\) 给出．证明数列 \(\left\{x_{n}\right\}_{n \geq 1}\)是单调数列，并求 \(\lim _{n \rightarrow \infty} x_{n}\) ．

\subsection{题目7}
7．（15 分）设 \(f: \mathbb{R}^{2} \rightarrow \mathbb{R}\) 在 \((0,0)\) 附近存在偏导数 \(\frac{\partial f}{\partial x}\) 和 \(\frac{\partial f}{\partial y}\) ．设 \(\frac{\partial f}{\partial x}\) 在 \((0,0)\) 点是连续的．求证：对于任何 \(u, v \in \mathbb{R},\left.\frac{\mathrm{~d}}{\mathrm{~d} t}\right|_{t=0} f(t u, t v)=u \frac{\partial f}{\partial x}(0,0)+v \frac{\partial f}{\partial y}(0,0)\) ．

\subsection{题目8}
8．（15 分）设 \(f(x)\) 在 \([-1,1]\) 上二次可导，并有 \(f(-1)=0, f(0)=0, f(1)=1\) 。证明：存在 \(\xi \in(-1,1)\) ，使得 \(f^{\prime \prime}(\xi)=1\) ．

\subsection{题目9}
9．（15 分）设 \(f(x)\) 是 \(2 \pi\) 为周期的黎曼可积函数，设 \(a_{n}, b_{n}\) 是它的 Fourier 系数。求证：对于任何的\(c, d \in[-\pi, \pi]\) ，都有

\[
\int_{c}^{d} f(x) \mathrm{d} x=\int_{c}^{d} \frac{a_{0}}{2} \mathrm{~d} t+\sum_{n=1}^{\infty} \int_{c}^{d}\left(a_{n} \cos n x+b_{n} \sin n x\right) \mathrm{d} x .
\]

\subsection{题目10}
10．（15 分）设有实数列 \(\left\{a_{n}\right\}_{n \geq 1}\) ，令

\[
S_{n}=a_{1}+a_{2}+\cdots+a_{n}, \sigma_{n}=\frac{S_{1}+S_{2}+\cdots+S_{n}}{n} .
\]

（1）证明：\(\left\{S_{n}\right\}\) 有极限推出 \(\left\{\sigma_{n}\right\}\) 有相同的极限．\\
（2）如果 \(\left\{\sigma_{n}\right\}\) 收敛，而且 \(a_{n}=o\left(\frac{1}{n}\right)\) ，则 \(\left\{S_{n}\right\}\) 收敛。
\end{document}
